\chapter{Related Work}

\setlength{\parindent}{1em}

This chapter situates MiPay within the relevant literature. The system draws on
four research streams: (i) automatic speech recognition, particularly
large-scale multilingual models; (ii) large language models and instruction
following; (iii) structured information extraction and constrained decoding; and
(iv) Arabic and code-switched natural-language processing. After reviewing the
foundational and recent work in each stream, the chapter examines a set of
directly comparable studies using a consistent four-part breakdown
(methodology, results, main features, limitations), and closes with a
comparative analysis that positions MiPay against them.

\section{Foundations and Research Streams}

\subsection{Automatic Speech Recognition}

Speech recognition progressed in three broad waves. The first replaced
Gaussian-mixture acoustic models with deep neural networks inside a
hidden-Markov-model framework, cutting error rates substantially on benchmark
tasks~\cite{hinton2012deep} and crystallizing in toolkits such as
Kaldi~\cite{povey2011kaldi}. The second wave introduced end-to-end neural
training, removing the hand-engineered pronunciation lexicons and alignment
stages. Connectionist Temporal Classification (CTC) made it possible to train
sequence models without frame-level alignment~\cite{graves2006ctc}, and
self-supervised pre-training on unlabeled audio, exemplified by
wav2vec\,2.0~\cite{baevski2020wav2vec}, sharply reduced the amount of labeled
speech needed to reach competitive accuracy.

The third wave---and the one MiPay builds upon---is large-scale weakly
supervised sequence-to-sequence modeling. Whisper~\cite{radford2023whisper}
trains an encoder--decoder transformer~\cite{vaswani2017attention} on 680{,}000
hours of multilingual, multitask audio scraped from the web. Rather than chasing
a single benchmark, Whisper optimizes for \emph{robustness}: it generalizes
across accents, recording conditions, and languages without dataset-specific
fine-tuning, and it performs language identification, transcription, and
translation within one model. This robustness, and its strong zero-shot Arabic
and English performance, are precisely the properties a bilingual,
real-world voice application needs. Multilingual speech corpora such as Common
Voice~\cite{ardila2020commonvoice} and Arabic-specific benchmarks such as the
MGB-2 broadcast challenge~\cite{ali2016mgb2} provide the evaluation context
against which such models' Arabic capabilities are understood.

\subsection{Large Language Models and Instruction Following}

The transformer architecture~\cite{vaswani2017attention} underpins both modern
ASR and modern language modeling. On the language side, two ideas converged to
produce today's general-purpose models. First, scaling autoregressive language
modeling produced \emph{emergent few-shot ability}: GPT-3 showed that a
sufficiently large model could perform new tasks from a handful of in-context
examples with no gradient updates~\cite{brown2020gpt3,radford2019gpt2}. Second,
\emph{instruction tuning} aligned these models with human task descriptions:
FLAN demonstrated that fine-tuning on a diverse set of instructions improves
zero-shot generalization~\cite{wei2022flan}, and reinforcement learning from
human feedback (InstructGPT) further aligned model behavior with user
intent~\cite{ouyang2022instructgpt}. Prompting strategies such as chain-of-thought
reasoning~\cite{wei2022cot} and the broader prompting
literature~\cite{liu2023promptsurvey} expanded what could be elicited without
retraining.

Open-weight models brought these capabilities on-premises. LLaMA showed that
carefully trained smaller models could rival much larger closed
ones~\cite{touvron2023llama}, and the Qwen series~\cite{yang2024qwen2,qwen2024qwen25}
pushed strong multilingual---including Arabic---instruction following down to the
7-billion-parameter scale that fits on commodity hardware. Efficient-inference
research makes such local deployment practical: 8-bit and 4-bit post-training
quantization~\cite{dettmers2022llmint8,frantar2023gptq} shrinks memory footprints
with little quality loss, which is what allows MiPay to host both a speech model
and a language model on a single 16\,GB machine.

\subsection{Structured Extraction and Constrained Decoding}

Turning free text into typed records is information extraction, classically
framed as named-entity recognition and slot filling. Neural sequence labelers
such as BiLSTM-CRF~\cite{lample2016ner} and the large body of deep-learning NER
work surveyed by Li et al.~\cite{li2022ner} achieve high accuracy but depend on
substantial labeled training data. The generative-LLM alternative trades that
data requirement for prompt engineering, but introduces a new failure mode:
free-form generation can violate the required output format.

Constrained (guided) decoding closes this gap. By intersecting the model's
token distribution at each step with a finite-state machine compiled from a
regular expression, grammar, or JSON schema, the decoder can only emit
continuations that keep the output valid~\cite{willard2023outlines}.
Grammar-constrained decoding has been shown to improve structured-NLP task
performance without any fine-tuning~\cite{geng2023grammar}. MiPay uses exactly
this mechanism---Ollama's schema-constrained \texttt{format} option---so that
extraction output is guaranteed to be valid JSON conforming to the transaction
schema, and enum fields can never take an illegal value.

\subsection{Arabic and Code-Switched NLP}

Arabic NLP is complicated by morphological richness, diglossia, and orthographic
variation. Tooling such as Farasa~\cite{abdelali2016farasa} and CAMeL
Tools~\cite{obeid2020camel} addresses segmentation and morphological analysis,
while AraBERT~\cite{antoun2020arabert} and multilingual encoders
XLM-R~\cite{conneau2020xlmr} provide contextual representations for Arabic
understanding tasks. Code-switching---fluidly mixing two languages---is itself a
distinct research area~\cite{sitaram2019codeswitch}, posing problems for systems
that assume a single language per input. MiPay's design responds to this
literature pragmatically: rather than building a dedicated code-switching model,
it leans on Whisper's and Qwen's multilingual pre-training to absorb mixed input,
and isolates the brittle, language-specific operations (Arabic-Indic numeral
mapping, dialectal date phrases) in a deterministic, separately tested layer.

\section{Summary of Recent Studies}

The following studies are the most directly comparable to MiPay. Each is reviewed
under a consistent breakdown to keep the later comparison meaningful.

\subsection{Robust Speech Recognition via Large-Scale Weak Supervision (Whisper)}

\subsubsection{Methodology}
Radford et al.~\cite{radford2023whisper} train an encoder--decoder transformer on
680{,}000 hours of weakly labeled multilingual audio collected from the web,
formulating transcription, translation, language identification, and
voice-activity detection as a single multitask sequence-prediction problem
conditioned on special task tokens.

\subsubsection{Results}
The models approach human-level robustness on out-of-distribution English
benchmarks in a zero-shot setting and deliver strong multilingual performance,
including competitive Arabic transcription, without per-dataset fine-tuning.
Accuracy scales monotonically with model size from \texttt{tiny} to
\texttt{large}.

\subsubsection{Main Features}
A single model covering 99 languages; robustness to noise and accent; built-in
language detection; permissive open release of weights.

\subsubsection{Limitations}
Inference is comparatively heavy for the larger sizes; the model can hallucinate
text during long silences (mitigated by voice-activity filtering); and it detects
a single dominant language per clip rather than explicitly modeling
code-switching.

\subsubsection{Implications and Contributions}
Whisper is the speech backbone of MiPay. Its open weights and multilingual
robustness are what make a self-hosted bilingual ASR stage feasible; its CPU cost
motivates the use of \texttt{faster-whisper}/CTranslate2 and quantization.

\subsection{Qwen2.5: An Open Multilingual Instruction-Tuned LLM}

\subsubsection{Methodology}
The Qwen2.5 technical report~\cite{qwen2024qwen25} describes a family of
decoder-only transformers pre-trained on a large multilingual corpus and then
instruction-tuned and preference-aligned, released openly across a range of
parameter scales including 3B and 7B.

\subsubsection{Results}
The models report strong results on multilingual understanding, reasoning, coding,
and structured-output benchmarks, with the 7B instruct variant offering a
favorable quality-to-resource ratio and solid Arabic instruction following.

\subsubsection{Main Features}
Open weights at deployable scales; strong multilingual coverage; reliable
instruction following and JSON formatting; compatibility with local servers such
as Ollama and with quantized inference.

\subsubsection{Limitations}
Like all LLMs it can produce semantically incorrect values, is unreliable at
exact arithmetic (including date computation), and its quality degrades at the
smaller 3B scale used as a low-resource fallback.

\subsubsection{Implications and Contributions}
Qwen2.5-7B-Instruct is MiPay's extraction engine. Its multilingual competence
removes the need for task-specific training data, while its known arithmetic
unreliability directly motivates MiPay's deterministic date-resolution layer.

\subsection{Efficient Guided Generation for Large Language Models}

\subsubsection{Methodology}
Willard and Louf~\cite{willard2023outlines} formulate constrained generation as
indexing a finite-state machine over the model's vocabulary, enabling
regex- and JSON-schema-guided decoding with negligible per-token overhead.

\subsubsection{Results}
The method guarantees that generated text matches the supplied
grammar/schema while preserving generation speed, eliminating post-hoc parsing
and retry loops.

\subsubsection{Main Features}
Output validity by construction; support for regular expressions and JSON
schemas; model-agnostic applicability.

\subsubsection{Limitations}
Constraints govern \emph{form}, not \emph{semantics}: a schema-valid output can
still be factually wrong, so a downstream validation or human-review step remains
necessary.

\subsubsection{Implications and Contributions}
This is the theoretical basis for MiPay's reliance on schema-constrained
decoding. It justifies the architectural decision to treat the LLM's JSON as
always parseable and to place semantic safeguards in post-processing and human
confirmation rather than in defensive parsing.

\subsection{Grammar-Constrained Decoding for Structured NLP Tasks}

\subsubsection{Methodology}
Geng et al.~\cite{geng2023grammar} apply context-free-grammar constraints during
decoding to a range of structured NLP tasks, without fine-tuning the underlying
model.

\subsubsection{Results}
Constrained decoding improves output validity and task accuracy relative to
unconstrained prompting on several structured benchmarks.

\subsubsection{Main Features}
No fine-tuning required; general grammar support; consistent gains on
structure-sensitive tasks.

\subsubsection{Limitations}
Designing a correct grammar/schema is itself non-trivial, and overly tight
constraints can suppress otherwise-correct outputs.

\subsubsection{Implications and Contributions}
Reinforces the schema-first extraction strategy and validates that constraint at
decode time, rather than correction after the fact, is the principled way to get
structured output from a general model.

\subsection{AraBERT and Arabic Language Understanding}

\subsubsection{Methodology}
Antoun et al.~\cite{antoun2020arabert} pre-train a BERT-style encoder on a large
Arabic corpus with Arabic-specific preprocessing (e.g., Farasa
segmentation~\cite{abdelali2016farasa}).

\subsubsection{Results}
AraBERT outperforms multilingual baselines on Arabic sentiment, NER, and
question-answering tasks, establishing the value of language-specific
pre-training and preprocessing.

\subsubsection{Main Features}
Strong Arabic representations; demonstrated benefit of morphological
preprocessing; open release.

\subsubsection{Limitations}
As an encoder it performs classification/labeling rather than generation, and it
requires labeled data for each downstream slot-filling task---precisely the data
MiPay does not have for bilingual financial extraction.

\subsubsection{Implications and Contributions}
AraBERT represents the supervised alternative MiPay deliberately defers to future
work. It also informs MiPay's Arabic normalization choices in the evaluation
harness (diacritic stripping, alef/ta-marbuta unification).

\section{Comparative Analysis of Related Studies}

Table~\ref{tab:related-comparison} compares the reviewed approaches along the
axes that matter for a self-hosted bilingual finance application: the task each
addresses, whether it needs task-specific labeled data, its native handling of
Arabic and code-switching, whether it can run self-hosted on commodity hardware,
and whether it yields guaranteed-structured output.

\begin{table}[H]
\centering
\caption{Comparison of related approaches against MiPay's requirements.}
\label{tab:related-comparison}
\small
\renewcommand{\arraystretch}{1.3}
\setlength{\tabcolsep}{4pt}
\begin{tabular}{
  >{\raggedright\arraybackslash}p{2.4cm}
  >{\raggedright\arraybackslash}p{1.7cm}
  >{\raggedright\arraybackslash}p{2.1cm}
  >{\centering\arraybackslash}p{1.6cm}
  >{\centering\arraybackslash}p{1.4cm}
  >{\centering\arraybackslash}p{1.9cm}
}
\toprule
\textbf{Approach} & \textbf{Primary task} & \textbf{Needs labeled data} & \textbf{Arabic / mixed} & \textbf{Self-host} & \textbf{Structured out} \\
\midrule
Whisper~\cite{radford2023whisper} & Speech-to-text & No (zero-shot) & Strong & Yes & N/A \\
Qwen2.5-7B~\cite{qwen2024qwen25} & General LLM & No (few-shot) & Good & Yes & Via prompting \\
Guided decoding~\cite{willard2023outlines} & Output control & No & N/A & Yes & Guaranteed \\
BiLSTM-CRF NER~\cite{lample2016ner} & Slot filling & Yes (large) & Needs corpus & Yes & Tags only \\
AraBERT~\cite{antoun2020arabert} & Arabic understanding & Yes & Strong (MSA) & Yes & Tags only \\
Rule/regex systems & Pattern match & No & Brittle & Yes & Rigid \\
\midrule
\textbf{MiPay (this work)} & \textbf{Voice $\rightarrow$ records} & \textbf{No} & \textbf{Native (mixed)} & \textbf{Yes} & \textbf{Guaranteed} \\
\bottomrule
\end{tabular}
\end{table}

The comparison clarifies MiPay's position. No single prior system delivers the
full requirement---bilingual voice in, validated structured financial records
out, self-hosted, without labeled training data. MiPay's contribution is the
\emph{composition}: Whisper supplies robust multilingual transcription; a
few-shot, schema-constrained Qwen2.5-7B supplies training-data-free structured
extraction with guaranteed-valid output; and a deterministic post-processing
layer supplies the exactness (numerals, dates, currency) that neither the speech
model nor the language model provides reliably. The supervised
approaches~\cite{lample2016ner,antoun2020arabert} are not competitors but a
documented future direction, to be pursued once user-corrected data accumulates
(Chapter~6). This synthesis---rather than any one component---is what the
remainder of the thesis designs, implements, and evaluates.